% 自動生成（tools/mkchecklist.py p3）。各章の「ふりかえり」の項目を集めたもの（第13回）
\begingroup\footnotesize
\def\bkcb{\tikz[baseline=-0.5ex]\draw[\bkTF, line width=0.7pt, rounded corners=0.5pt] (0,-0.3em) rectangle (0.8em,0.5em);}
\begin{longtblr}{colspec={Q[c,m,wd=11mm] X[l,m] Q[c,m,wd=10mm] Q[c,m,wd=9mm] Q[c,m,wd=9mm]}, width=\linewidth, rowhead=1, colsep=3pt,
  row{1}={bg=bkVol!14, font=\sffamily\bfseries\scriptsize}, hlines={0.4pt,bkLine}, rowsep=0.7pt}
章 & できるようになったこと & ページ & 1回め & 2回め \\
{\sffamily\bfseries\color{\bkTT}波-0} & ゆれの周期と振動数、波の波長と速さを、身近な例で説明できる。 & \pageref{ph:波-0:sec2} & \bkcb & \bkcb \\
 & 何を無視したかを書いてから、波が届くまでの時間や距離を、7つの段で求められる。 & \pageref{ph:波-0:ex3} & \bkcb & \bkcb \\
 & 身近な「なぜ？」の答えが、この巻のどの章にあるかを探せる。 & \pageref{ph:波-0:sec6} & \bkcb & \bkcb \\
{\sffamily\bfseries\color{\bkTT}波-1} & ふりこの1往復の時間が、ふりこの長さで決まり、おもりの重さにはよらないことを、実験で確かめられる。 & \pageref{ph:波-1:sec3} & \bkcb & \bkcb \\
 & ばねとおもり・ふりこの運動方程式から、単振動の式と周期を導き、エネルギーで速さを求められる。 & \pageref{ph:波-1:sec2} & \bkcb & \bkcb \\
 & 等価ばね定数・静たわみ・実体ふりこ・ねじり振動で、いろいろな物の固有振動数を見積もれる。 & \pageref{ph:波-1:sec4} & \bkcb & \bkcb \\
 & ゆれの記録から減衰比を読み、力学・電気・熱・流体の相似の表で、同じ形の式を見分けられる。 & \pageref{ph:波-1:sec7} & \bkcb & \bkcb \\
{\sffamily\bfseries\color{\bkTT}波-2} & 共振（共鳴）を、ゆする速さとゆれの大きさの関係で説明できる。 & \pageref{ph:波-2:sec1} & \bkcb & \bkcb \\
 & 表計算でばねをゆすって、共振曲線とゆれのおくれを読みとれる。 & \pageref{ph:波-2:sec3} & \bkcb & \bkcb \\
 & 回る物の不つり合い・でこぼこ道・風の渦が、どれだけの力で、1秒に何回ゆするかを求められる。 & \pageref{ph:波-2:kazoeru} & \bkcb & \bkcb \\
 & 共振の倍率・Q 値・伝達率を求め、防振ばねを選べる。 & \pageref{ph:波-2:sec5} & \bkcb & \bkcb \\
 & 不つり合い・回る軸・道のでこぼこ・風の渦によるゆれを見積もり、ボード線図で読める。 & \pageref{ph:波-2:sec7} & \bkcb & \bkcb \\
{\sffamily\bfseries\color{\bkTT}波-3} & ゆれが行き来するわけを、2つのモードのたし算で説明できる。 & \pageref{ph:波-3:sec2} & \bkcb & \bkcb \\
 & 表計算で2自由度のゆれと、動吸振器の効きめを確かめられる。 & \pageref{ph:波-3:sec3} & \bkcb & \bkcb \\
 & ラグランジュの式で、エネルギーから運動方程式を立てられる。 & \pageref{ph:波-3:sec4} & \bkcb & \bkcb \\
 & 質量行列のある固有値問題で、モードを求められる。 & \pageref{ph:波-3:sec5} & \bkcb & \bkcb \\
 & 動吸振器の同調と、モードの直交・モード解析を説明できる。 & \pageref{ph:波-3:sec6} & \bkcb & \bkcb \\
{\sffamily\bfseries\color{\bkTT}波-4} & フィードバックのしくみを、ブロック線図で説明できる。 & \pageref{ph:波-4:sec2} & \bkcb & \bkcb \\
 & 表計算で部屋の温度の比例・積分の制御を追える。 & \pageref{ph:波-4:sec3} & \bkcb & \bkcb \\
 & 物理の法則から伝達関数を作り、応答の速さを読める。 & \pageref{ph:波-4:sec4} & \bkcb & \bkcb \\
 & 残るずれと PID の働きを、最終値の定理で説明できる。 & \pageref{ph:波-4:sec6} & \bkcb & \bkcb \\
 & ラウスの表で、安定かどうかを判定できる。 & \pageref{ph:波-4:sec8} & \bkcb & \bkcb \\
 & ゲイン余裕・位相余裕と状態フィードバックで、安定の度合いを考えられる。 & \pageref{ph:波-4:sec9} & \bkcb & \bkcb \\
{\sffamily\bfseries\color{\bkTT}波-5} & 波で進むのは「ゆれの形」だと、横波・縦波で説明できる。 & \pageref{ph:波-5:sec1} & \bkcb & \bkcb \\
 & 振幅・波長・周期・振動数と $c=f\lambda$ を使いこなせる。 & \pageref{ph:波-5:sec2} & \bkcb & \bkcb \\
 & 正弦波の式 $y=A\sin(\omega t-kx)$ からゆれを求められる。 & \pageref{ph:波-5:sec3} & \bkcb & \bkcb \\
 & 媒質の運動方程式から波動方程式と波の速さを導ける。 & \pageref{ph:波-5:sec4} & \bkcb & \bkcb \\
 & おもりの列の分散関係・群速度と、計算の刻みの条件がわかる。 & \pageref{ph:波-5:sec7} & \bkcb & \bkcb \\
{\sffamily\bfseries\color{\bkTT}波-6} & 固定端・自由端の反射と、定在波の腹と節を説明できる。 & \pageref{ph:波-6:sec1} & \bkcb & \bkcb \\
 & 弦の固有振動数を求め、フレットと音の高さの関係がわかる。 & \pageref{ph:波-6:sec3} & \bkcb & \bkcb \\
 & 閉管・開管の固有振動数と共鳴から、音の速さを求められる。 & \pageref{ph:波-6:sec4} & \bkcb & \bkcb \\
 & 変数分離とフーリエ級数で、はじいた弦の倍音を求められる。 & \pageref{ph:波-6:sec5} & \bkcb & \bkcb \\
 & はりの振動の式を解き、鉄琴や配管の固有振動数を求められる。 & \pageref{ph:波-6:sec9} & \bkcb & \bkcb \\
{\sffamily\bfseries\color{\bkTT}波-7} & 波が運ぶエネルギーと強さを、振幅・振動数・距離から求められる。 & \pageref{ph:波-7:sec1} & \bkcb & \bkcb \\
 & インピーダンスの比から、境目ではね返る割合と通りぬける割合を求められる。 & \pageref{ph:波-7:sec4} & \bkcb & \bkcb \\
 & 表計算で、2本の弦をつないだときの山のはね返りと通りぬけを追える。 & \pageref{ph:波-7:sec4} & \bkcb & \bkcb \\
 & 波の式と境界条件から反射係数を導き、固定端・自由端・はね返らない端を1つの式で説明できる。 & \pageref{ph:波-7:sec5} & \bkcb & \bkcb \\
 & 吸収による弱まり、整合層、光の反射率を、インピーダンスと dB で見積もれる。 & \pageref{ph:波-7:sec7} & \bkcb & \bkcb \\
{\sffamily\bfseries\color{\bkTT}波-8} & 近い高さの2つの音が重なるときの、うなりの回数を求められる。 & \pageref{ph:波-8:sec2} & \bkcb & \bkcb \\
 & ホイヘンスの原理で、反射・屈折の法則と回折を説明できる。 & \pageref{ph:波-8:sec1} & \bkcb & \bkcb \\
 & 2つの波源からの道のりの差で、強め合う所と弱め合う所を求められる。 & \pageref{ph:波-8:sec3} & \bkcb & \bkcb \\
 & 波面を数えてドップラー効果の式を導き、聞こえる高さを求められる。 & \pageref{ph:波-8:sec5} & \bkcb & \bkcb \\
 & 斜めのドップラー効果・反射板・衝撃波で、速さをはかるしくみを説明できる。 & \pageref{ph:波-8:sec7} & \bkcb & \bkcb \\
{\sffamily\bfseries\color{\bkTT}波-9} & 雷やこだまの時間から、音の速さを使って距離を求められる。 & \pageref{ph:波-9:sec2} & \bkcb & \bkcb \\
 & 音の大きさを dB で表し、音を合わせたときや、離れたときの大きさを計算できる。 & \pageref{ph:波-9:sec7} & \bkcb & \bkcb \\
 & 縮みにくさと密度から気体・液体・固体の音の速さを求め、遮音と残響を見積もれる。 & \pageref{ph:波-9:sec3} & \bkcb & \bkcb \\
{\sffamily\bfseries\color{\bkTT}波-10} & こだまの時間から、コウモリ・ソナー・超音波探傷の距離を求められる。 & \pageref{ph:波-10:sec1} & \bkcb & \bkcb \\
 & 振動の変位・速度・加速度を換算し、回転機械のスペクトルから故障の見当をつけられる。 & \pageref{ph:波-10:sec6} & \bkcb & \bkcb \\
 & 1自由度の模型で、加速度計と変位計のしくみを説明し、感度のずれを計算できる。 & \pageref{ph:波-10:sec7} & \bkcb & \bkcb \\
{\sffamily\bfseries\color{\bkTT}波-11} & 音をサンプリングと量子化で数の列にし、データの量を求められる。 & \pageref{ph:波-11:sec1} & \bkcb & \bkcb \\
 & エイリアシングで、車輪が逆に回って見えるわけを説明できる。 & \pageref{ph:波-11:sec2} & \bkcb & \bkcb \\
 & 目もり $\Delta f=f_{\mathrm s}/N$ を読み、窓関数でもれを小さくできる。 & \pageref{ph:波-11:sec4} & \bkcb & \bkcb \\
 & サンプリング定理・フィルター・平均で、S/N 比を見積もれる。 & \pageref{ph:波-11:sec5} & \bkcb & \bkcb \\
 & 情報量と通信路容量を計算し、熱のエントロピーと比べられる。 & \pageref{ph:波-11:sec9} & \bkcb & \bkcb \\
{\sffamily\bfseries\color{\bkTT}波-12} & 光の直進・反射・屈折を、身近な例と作図で説明できる。 & \pageref{ph:波-12:sec1} & \bkcb & \bkcb \\
 & 屈折の法則で、屈折角・臨界角・見かけの深さを求められる。 & \pageref{ph:波-12:sec3} & \bkcb & \bkcb \\
 & レンズの式で、像の位置と倍率を求め、光学機器の倍率を見積もれる。 & \pageref{ph:波-12:sec6} & \bkcb & \bkcb \\
 & 屈折率が色で変わること（分散）から、プリズムと虹の色を説明できる。 & \pageref{ph:波-12:sec8} & \bkcb & \bkcb \\
 & フェルマーの原理から反射・屈折の法則を導き、光線の行列でレンズの組を計算できる。 & \pageref{ph:波-12:sec7} & \bkcb & \bkcb \\
{\sffamily\bfseries\color{\bkTT}波-13} & ヤングの実験・回折格子・薄い膜で、光が強め合う所と弱め合う所を、経路の差から計算できる。 & \pageref{ph:波-13:sec2} & \bkcb & \bkcb \\
 & 円い穴の回折で見分けられる角の限界と、偏光板を通る光の強さ・ブルースター角を計算できる。 & \pageref{ph:波-13:sec5} & \bkcb & \bkcb \\
 & ゆすられる電子のモデル（強制振動）から、空が青く夕焼けが赤いわけを説明できる。 & \pageref{ph:波-13:sec8} & \bkcb & \bkcb \\
{\sffamily\bfseries\color{\bkTT}波-14} & 電波の振動数からアンテナの大きさを、電波が往復する時間から距離を、見積もれる。 & \pageref{ph:波-14:sec5} & \bkcb & \bkcb \\
 & AM・FM・ディジタルの変調を、波の振幅・振動数・位相を変えることとして説明できる。 & \pageref{ph:波-14:sec2} & \bkcb & \bkcb \\
 & たくさんの周波数で同時に送る方式（OFDM）の速さを計算し、通信路容量の上限とくらべられる。 & \pageref{ph:波-14:ex3} & \bkcb & \bkcb \\
\end{longtblr}
\endgroup
